\section{Algorithms and data structures}

\subsection{What is an algorithm?}
An \textbf{algorithm} is a sequence of instructions to perform a computation. 
Algorithms can be expressed in several different ways such as flowcharts, pseudo-code, and working code snippets. The sequence of operations must be expressed unambiguously. Using the right algorithm sometimes makes all the difference between being able to attack a problem in a practical time or not.

\subsubsection{Example: sequential vs. binary search}
Let's consider a classic problem: finding an element in a sorted list. Let's try and solve it in the two most intuitive ways.

The first method is the sequential search which utilizes a bit of brute force: you loop over the list until you find (or do not find) the target.
    
The second method is the Binary search, schematically represented below:
\begin{itemize}
    \item[-] Start from the middle (if that's the target you're done)
    \item[-] If the target is smaller (larger) than the element in the center, bisect the half-list on the left (right)
    \item[-] Iterate until you've found the target (or exhausted the list)
\end{itemize}
\begin{center}
    \includegraphics[width=0.6\textwidth]{binary_search}
\end{center}
As an exercise, try to determine which algorithm is more efficient. if you think about it, this problem is essentially similar to that of finding a person in the phone book.

\subsection{Complexity of an algorithm}

As an example, let's try to find the largest element in a list of length $n$ implementing this problem in Python:

\input{pygments/find_max}

How many fundamental instructions does this code execute? You should realize that this is an ill-posed question, since the answer depends on the structure of the list. Nonetheless, let's count:
\begin{itemize}
    \item[-] One assignment and one list lookup at the beginning: $2$
    \item[-] One lookup, one assignment and one comparison for each iteration in the for loop: $3(n - 1)$
    \item[-] A variable number of assignments: between $0$ and $(n - 1)$
    \item[-] One final return instruction: 1
\end{itemize}
Answer: anything between $3n$ and $4n - 1$, depending on the input list.

What does this example tell us?

\underline{Message \#1:} \textit{The exact number of fundamental instructions that an algorithm performs is not determined a priori}. It depends on the input data, instead, and so does the running time.

There are a few questions that you can legitimately ask, anyway:
\begin{itemize}
    \item[-] How many instructions in the \textit{worst case}?
    \item[-] How many instructions in the \textit{best case}?
    \item[-] How many instructions on \textit{average}?
\end{itemize}

\underline{Message \#2:} \textit{The exact number of operations doesn't really matter}, does it? Different machines have different executions speeds and different languages have different meanings of \textit{fundamental instruction}.

\underline{Message \#3:} Still, \textit{the running time is related to the number of fundamental instructions}.

With these messages in mind, how can we define how "fast" an algorithm is? The easiest way is to analyze its asymptotic behavior.

\subsubsection{Asymptotic behavior and big-O notation}

Say you have an algorithm operating on an input of length $n$ (e.g., a list with $n$ elements or a string with $n$ characters).
How many fundamental instructions $N$ does it take for your algorithm to run? \footnote{As said before, this question is ambiguous. The correct one should be "How many fundamental instructions $N$ does it take for your algorithm to run \textit{in the worst case}?" (easiest to compute) or "\textit{in the average case}?".} When asking this question, you are trying to find a function $f$ such that:
\begin{equation*}
    N = f(n)    
\end{equation*}
Since the real problems in algorithm implementation come around when $n$ is really high, we are interested in the \textbf{asymptotic behaviour} of f: drop all the terms that grow slowly with $n$ and only keep the one that grows faster:
\begin{equation*}
    4n - 1 \approx 4n \quad\text{and}\quad 2n^2 + 6n + 3 \approx 2n^2 \quad \text{for} \quad n \xrightarrow{} +\infty
\end{equation*}
Let's go one step further, and say that we neglect the multiplicative factor in front of the leading term:
\begin{equation*}
    4n - 1 \approx n \quad\text{and}\quad 2n^2 + 6n + 3 \approx n^2    
\end{equation*}
That's the so-called \textbf{big-O notation}: the two algorithms in the example are $O(n)$ (the less complex one) and $O(n^2)$ (the more complex one). Your goal should be to write an algorithm with the lowest possible complexity. Take a look at the graph below: if $n = 10^6$ and you can beat down the complexity from $n^2$ to $n\log(n)$ you are cutting down the execution time by one million.

\begin{center}
    \includegraphics[width=0.5\textwidth]{bigo}
\end{center}

Now, how can you measure the asymptotic behaviour of an algorithm?
\begin{itemize}
    \item[-] By brute force: implement the algorithm, run it on input data of different sizes and time the run (be careful: results may vary from run to run) and plot the running time vs. input size.
    \item[-] By analysis: count the instructions and evaluate the best, worst and average case (this can be difficult for complex programs, and subject to the idiosyncrasies of the language).
    \item[-] "By eye" (e.g. one loop: $O(n)$; two sequential loops: $O(n)$; two nested loops: $O(n^2)$)
\end{itemize}
Here's a fun article about the complexity of songs: \url{http://www.cs.bme.hu/~friedl/alg/knuth_song_complexity.pdf}

\subsection{Data structures}
During this course you will have to work on real-life data. In which "objects" can data be stored? How fast can you operate on such objects ($\equiv$ how complex are the algorithms that let you interact with such data)? Here are two examples of \textbf{data structures} that you should master while coding in Python: lists and hash tables (dictionaries).

\subsubsection{Lists}

The following table summarizes the time complexity of common operations on Python lists:

\begin{center}
\begin{tabular}{lll}
    \hline
    Operation & Average case & Worst case\\
    \hline
    \hline
    Copy & O(n) & O(n)\\
    Append & O(1) & O(1)\\
    Insert & O(n) & O(n)\\
    Get Item & O(1) & O(1)\\
    Set Item & O(1) & O(1)\\
    Delete Item & O(n) & O(n)\\
    Iteration & O(n) & O(n)\\
    min(s), max(s) & O(n) & \\
    Get Length & O(1) & O(1)\\
    \hline
\end{tabular}
\end{center}

 Do you remember how to index a list? You should, but here is a simple figure:
 
\begin{center}
    \includegraphics[width=0.5\textwidth]{python-list}
\end{center}

\subsubsection{Hash tables}
A \textbf{hash table} is a data structure that implements an associative array, a structure that can map keys to values. It is designed to optimize data retrieval, allowing for extremely fast access to data regardless of the size of the dataset.

\begin{center}
    \includegraphics[width=0.5\textwidth]{hash_table}
\end{center}

The efficiency of a hash table relies on a hashing function. When you insert a key-value pair (e.g., \texttt{\{'name': 'Alice'\}}), the system performs the following steps:
\begin{enumerate}
    \item[-] Hashing: The key ('name') is passed through a hash function, which converts it into a unique integer (the hash).
    \item[-] Indexing: This integer is used to calculate an index (usually via the modulo operator) corresponding to a specific slot, or bucket, in an array.
    \item[-] Storage: The value ('Alice') is stored in that specific bucket.
\end{enumerate}

When you want to retrieve the value, the system re-calculates the hash of the key, finds the index, and instantly retrieves the value without searching through the entire list.

Why do I even bring this up? Because Python dictionaries are highly optimized hash tables. The performance of dictionary operations is summarized in the following table:

%\begin{center}
%    \includegraphics[width=0.5\textwidth]{python-dict}
%\end{center}

\begin{center}
\begin{tabular}{lll}
    \hline
    Operation & Average case & Worst case\\
    \hline
    \hline
    Copy & O(n) & O(n)\\
    Get Item & O(1) & O(n)\\
    Set Item & O(1) & O(n)\\
    Delete Item & O(1) & O(n)\\
    Iteration & O(n) & O(n)\\
    \hline
\end{tabular}
\end{center}

Since the array size is finite and the set of possible keys is infinite, a \textbf{collision} occurs when two distinct keys result in the same index. Python can solve this problem thanks to a method called \textit{open addressing} (for more info about open addressing visit \url{https://adamgold.github.io/posts/python-hash-tables-under-the-hood/}).

Note that when a float corresponds to an integer, its hash is the same as that of the integer. The hash is used to map keys into indices, therefore $3$ and $3.$ are the same key to a dictionary. For example:

\input{pygments/dict_hashing}

\bigskip

\subsection{A real-life example: sorting}

Let's consider a naive implementation of a sorting algorithm:

\input{pygments/sloppy_sort}

Say you are asked to implement an algorithm to sort a list. Chances are that you would come up with two nested loops, O($n^2$). However, there is already a rich variety of sorting algorithms on the market, so you should focus on mastering them instead of creating one of your own.

The following table summarizes the time and space complexity of several popular sorting algorithms:

\begin{center}
    \includegraphics[width=\textwidth]{sorting_algs}
\end{center}

The \texttt{.sort()} method for lists in Python uses a hybrid algorithm derived from merge sort and insertion sort called \textbf{timsort} \footnote{The name comes from Tim Peters.}. It works by finding subsequences of the data that are already ordered and using that knowledge to sort the remainder more efficiently (\textit{"Although practicality beats purity"}, The Zen of Python).

Here is a quick example of how to use the built-in sorting method in Python:

\input{pygments/timsort}

\subsection{References}
\begin{itemize}
    \item[-] \url{https://en.wikipedia.org/wiki/Algorithm}
    \item[-] \url{https://discrete.gr/complexity/}
    \item[-] \url{https://wiki.python.org/moin/TimeComplexity}
    \item[-] \url{https://en.wikipedia.org/wiki/Sorting_algorithm}
    \item[-] \url{https://bugs.python.org/file4451/timsort.txt}
    \item[-] \url{https://en.wikipedia.org/wiki/Timsort}
\end{itemize}